%%%%%%%%%%  scan idx 217  |  folio 204  %%%%%%%%%%
\[
  i_s^*(P^n_{X/S})\simeq P^n_{X_s},
\]
d'où un
homomorphisme de restriction
\[
  i_s^*:\Diff^n(E,F)
  \longrightarrow
  \Diff^n\bigl(i_s^*(E),i_s^*(F)\bigr),
\]
pour $E,F$ des $\cO_X$-Modules localement libres de type fini (et un homomorphisme analogue pour les opérateurs à coefficients complexes). Soit $E$ un complexe, à degrés bornés, d'opérateurs différentiels à coefficients complexes sur $X$; pour que $E$ soit elliptique (1.2.5) il faut et il suffit que, pour tout $s\in S$, $i_s^*(E)$ soit un complexe elliptique (vérification laissée au lecteur).

\subsection*{2. Lemmes de rigidité}

\textbf{2.0.} Il est «~bien connu~» que les variétés différentiables compactes sont «~rigides~», i.e. n'admettent pas de déformation continue autre que triviale. Toutefois, comme ce résultat nous servira de façon essentielle au §3, et que les références semblent assez rares, il n'est peut-être pas inutile de donner des énoncés précis et quelques démonstrations.

On se fixe un espace topologique $S$. Les variétes mixtes dont il sera question seront supposées de classe $C^r$, avec $r\ge1$ (éventuellement $r=\infty$).

\textbf{2.1. Notations.} Soient $X,Y$ des $S$-variétés mixtes. On note $\Homop(X,Y)$ l'ensemble des $S$-morphismes de $X$ dans $Y$, $\underline{\Homop}(X,Y)$ le faisceau sur $X$ des germes de $S$-morphismes de $X$ dans $Y$ (défini par
\[
  \underline{\Homop}(X,Y)(U)=\Homop(U,Y)
\]
pour $U$ ouvert de $X$).

Soit $s\in S$. On note $i_s$ l'inclusion de $X_s$ dans $X$. On a un homomorphisme canonique d'induction
\begin{equation}\tag{2.1.1}
  i_s^{-1}\underline{\Homop}(X,Y)
  \longrightarrow
  \underline{\Homop}(X_s,Y_s),
\end{equation}
